\section{References}

\subsection*{Scripture, doctrine, law, and liturgy}
{\footnotesize
\begin{itemize}[itemsep=0.18em,topsep=0.25em]
\item Holy Scripture: Gen. 1--3; 6--9; 14:18--20; 17; 22; Exod. 12--14; 16--17; 19; 24; 28--30; Lev. 8--9; Num. 11; Deut. 6; 24; Josh. 3; 2~Kings 5; Pss. 22(23), 31(32), 50(51), 109(110), 127(128); Prov. 9; Song of Songs; Isa. 6; 11; 53--55; 61--62; Jer. 2--3; Ezek. 16; 18; 36--37; Hos. 1--3; Joel 2--3; Mal. 1:11; Tob. 7--8; 12; Matt. 3; 9; 16; 18--19; 22; 26; 28; Mark 3; 6; 10; 14; Luke 3; 7; 15; 22--24; John 2--3; 6; 10; 13--21; Acts 1--2; 6; 8--10; 13--14; 19--20; Rom. 5--8; 12; 1~Cor. 6--7; 10--12; 15; 2~Cor. 1; 4--5; 11; Gal. 3--4; Eph. 1--5; Col. 1--3; 1~Tim. 3--4; 2~Tim. 1; Titus 1; Heb. 2--10; Jas. 5; 1~Pet. 2--3; 2~Pet. 1; 1~John 1--2; Rev. 19; 21--22.
\item \emph{Catechism of the Catholic Church}, nos.~1076--1209, 1210--1666, especially 1113--1134 and 1210--1212, \sourceurl{https://www.vatican.va/content/catechism/en/part_two.html}{Holy See English text}. The exact English formula witnesses used here are \sourceurl{https://www.vatican.va/content/catechism/en/part_two/section_two/chapter_one/article_1/iii_how_is_the_sacrament_of_baptism_celebrated.html}{no.~1240} for Latin and Eastern Baptism, \sourceurl{https://www.vatican.va/content/catechism/en/part_two/section_two/chapter_one/article_2/ii_the_signs_and_the_rite_of_confirmation.html}{no.~1300} for current Latin and Byzantine Confirmation, and \sourceurl{https://www.vatican.va/content/catechism/en/part_two/section_two/chapter_two/article_4/vi_the_sacrament_of_penance_and_reconciliation.html}{no.~1449} for Latin absolution.
\item Council of Florence, Bull \emph{Exsultate Deo} / Decree for the Armenians (1439), DS 1310--1328; received as a major Latin doctrinal synthesis while later determinations govern the essential matter of Orders.
\item Council of Trent: sess.~V, Original Sin (1546), DS 1510--1516; sess.~VI, Justification, especially ch.~7 on God's glory and eternal life (1547), DS 1528--1529; sess.~VII, Decree on the Sacraments (1547), DS 1600--1613; sess.~XIII, Eucharist (1551), DS 1635--1661; sess.~XIV, Penance and Extreme Unction (1551), DS 1667--1719; sess.~XXI--XXII, Communion and the Sacrifice of the Mass (1562), DS 1725--1759; sess.~XXIII, Orders (1563), DS 1763--1778; sess.~XXIV, Matrimony (1563), DS 1797--1816.
\item Vatican Council II, \emph{Sacrosanctum Concilium} 5--7, 10, 26--27, 59--82, \sourceurl{https://www.vatican.va/archive/hist_councils/ii_vatican_council/documents/vat-ii_const_19631204_sacrosanctum-concilium_en.html}{text}; \emph{Lumen gentium} 7--11, 18--29, 35, 41, \sourceurl{https://www.vatican.va/archive/hist_councils/ii_vatican_council/documents/vat-ii_const_19641121_lumen-gentium_en.html}{text}; \emph{Orientalium Ecclesiarum} 2--6, 12--18, \sourceurl{https://www.vatican.va/archive/hist_councils/ii_vatican_council/documents/vat-ii_decree_19641121_orientalium-ecclesiarum_en.html}{text}; \emph{Presbyterorum ordinis} 2--6; \emph{Gaudium et spes} 47--52.
\item \emph{Code of Canon Law} (1983), cann.~834--1165, especially 840--846, 849, 861, 879--883, 897--911, 959--987, 998--1007, 1008--1024, 1055--1057, 1084, 1095--1108, 1112, 1116, 1127, 1141--1142, \sourceurl{https://www.vatican.va/archive/cod-iuris-canonici/cic_index_en.html}{Holy See English text}.
\item \emph{Code of Canons of the Eastern Churches} (1990), cann.~667--775 and 776--866, especially 675--697, 709--710, 722--729, 737--742, 741, 743--754, 776, 801, 825, 828, 832, 834, \sourceurl{https://www.vatican.va/content/john-paul-ii/la/apost_constitutions/documents/hf_jp-ii_apc_19901018_codex-can-eccl-orient-2.html}{official Latin text}.
\item Congregation for the Eastern Churches, \emph{Instruction for Applying the Liturgical Prescriptions of the Code of Canons of the Eastern Churches} (1996), nos.~42--51, 75--82, especially the unity of initiation at 50--51, \sourceurl{https://www.vatican.va/roman_curia/congregations/orientchurch/Istruzione/pdf/istruzione_inglese.pdf}{English PDF}.
\item Pius XII, apostolic constitution \emph{Sacramentum Ordinis} (1947), nos.~3--7, \sourceurl{https://www.vatican.va/content/pius-xii/la/apost_constitutions/documents/hf_p-xii_apc_19471130_sacramentum-ordinis.html}{official Latin text}.
\item Paul VI, apostolic constitution \emph{Divinae consortium naturae} (1971), essential rite of Confirmation, \sourceurl{https://www.vatican.va/content/paul-vi/la/apost_constitutions/documents/hf_p-vi_apc_19710815_divina-consortium.html}{official Latin text}; apostolic constitution \emph{Sacram unctionem infirmorum} (1972), \sourceurl{https://www.vatican.va/content/paul-vi/en/apost_constitutions/documents/hf_p-vi_apc_19721130_sacram-unctionem.html}{English text}.
\item Congregation for the Doctrine of the Faith, \emph{Note on the Minister of the Sacrament of the Anointing of the Sick} (2005), doctrinal note and commentary, \sourceurl{https://www.vatican.va/roman_curia/congregations/cfaith/documents/rc_con_cfaith_doc_20050211_unzione-infermi_en.html}{English text}; bishop or presbyter alone validly ministers the sacrament.
\item Dicastery for the Doctrine of the Faith, Note \emph{Gestis verbisque} on the validity of the sacraments (2024), nos.~8--27, \sourceurl{https://www.vatican.va/roman_curia/congregations/cfaith/documents/rc_ddf_doc_20240202_gestis-verbisque_en.html}{English text}; Congregation for the Doctrine of the Faith, Responses and Doctrinal Note concerning the formula of Baptism (2020), \sourceurl{https://www.vatican.va/roman_curia/congregations/cfaith/documents/rc_cdf_doc_20200624_responsum-nota-battesimo_en.html}{text}.
\item USCCB Secretariat of Divine Worship, \emph{Newsletter}, vol.~LVI (July--August 2020), 23, ``Five Questions about the CDF Response on the Baptismal Formula,'' \sourceurl{https://www.usccb.org/resources/newsletter-2020-07-and-08.pdf}{official conference PDF}: exact English conditional formula. This conference clarification is not represented as the universal liturgical book.
\item Joseph Deharbe, S.J., \emph{A Catechism of the Catholic Religion}, English trans., translator unnamed on the title page, ed.~P.~N. Lynch (New York: Schwartz, Kirwin \& Fauss, 1878), no.~556, p.~140, \sourceurl{https://archive.org/details/acatechismofcath00dehauoft}{edition scan}: exact historical English used for the unchanged 1962 Latin Confirmation form. The 1878 English is a public-domain human translation, not a project or AI translation and not a modern official English edition of the 1962 \emph{Pontificale Romanum}.
\item Pontifical Council for Promoting Christian Unity, \emph{Guidelines for Admission to the Eucharist between the Chaldean Church and the Assyrian Church of the East} (2001), on the Anaphora of Addai and Mari, \sourceurl{https://press.vatican.va/content/salastampa/it/bollettino/pubblico/2001/10/25/0586/01717.html}{text}.
\item \emph{Missale Romanum}, editio typica (1962), \emph{Missae votivae}, no.~11, \emph{Pro sponsis}, pp.~[75]--[77]: Collect, Secret, prayers after the Pater, Postcommunion, final blessing, and admonition; \sourceurl{https://media.churchmusicassociation.org/pdf/missale62.pdf}{facsimile}.
\item John Paul II, apostolic letter \emph{Ordinatio sacerdotalis} (1994), nos.~1--4, \sourceurl{https://www.vatican.va/content/john-paul-ii/en/apost_letters/1994/documents/hf_jp-ii_apl_19940522_ordinatio-sacerdotalis.html}{text}; apostolic exhortation \emph{Familiaris consortio} (1981), nos.~11--17, 49--64.
\item Leo XIII, encyclical \emph{Arcanum divinae sapientiae} (1880); 1917 \emph{Code of Canon Law}, can.~1013 \S1; Pius XI, encyclical \emph{Casti connubii} (1930), especially nos.~17--24 on the primary end of procreation and education and the ordered secondary goods, \sourceurl{https://www.vatican.va/content/pius-xi/en/encyclicals/documents/hf_p-xi_enc_19301231_casti-connubii.html}{Holy See English text}; Paul VI, encyclical \emph{Humanae vitae} (1968), nos.~8--14, \sourceurl{https://www.vatican.va/content/paul-vi/en/encyclicals/documents/hf_p-vi_enc_25071968_humanae-vitae.html}{text}.
\end{itemize}
}

\subsection*{Philosophical and scholastic framework}
{\footnotesize
\begin{itemize}[itemsep=0.18em,topsep=0.25em]
\item Aristotle, \emph{Physics} II.3 (the four causes); \emph{De anima} II.1 (act, potency, matter, and form in the account of soul); \emph{Metaphysics} V.2, VII, IX (cause, substance, act, and potency); \emph{Nicomachean Ethics} II, VIII.12 (habit, virtue, and spousal friendship); \emph{Politics} I.2 (household and natural association). Used selectively under Christian doctrinal correction.
\item Hugh of Saint Victor, \emph{De sacramentis christianae fidei} I.9; Peter Lombard, \emph{Sentences} IV, distinctions 1--42, as the scholastic bridge from patristic sign-language to the mature sacramental questions.
\item St. Thomas Aquinas, \emph{Summa theologiae} I, q.~4, a.~3 (participation), q.~43 (divine missions and indwelling); I--II, qq.~49--55 (habits and virtues), 62 (theological virtues), 109--114 (grace); III, q.~23 (adoptive filiation), qq.~60--65 (sacraments generally), 66--71 (Baptism), 72 (Confirmation), 73--83 (Eucharist), 84--90 (Penance), \sourceurl{https://www.newadvent.org/summa/4.htm}{working English index}.
\item Aquinas / posthumous \emph{Supplementum}, qq.~29--33 (Anointing), 34--40 (Orders), 41--68 (Matrimony), compiled principally from Aquinas's earlier \emph{Scriptum super Sententiis} IV; \emph{Summa contra Gentiles} IV.56--78; \emph{Commentary on Ephesians} 5, lects.~8--10.
\end{itemize}
}

\subsection*{Patristic foundations and the sacraments of initiation}
{\footnotesize
\begin{itemize}[itemsep=0.18em,topsep=0.25em]
\item \emph{Didache} 7; 9--10; 14; St. Justin Martyr, \emph{First Apology} 61, 65--67; St. Irenaeus, \emph{Against Heresies} III.17.1--3; IV.18.4--5; V.2.2--3.
\item Tertullian, \emph{De baptismo} 1--8; St. Cyprian of Carthage, \emph{Epistles} 69--75 (the rebaptism controversy, read with the Church's subsequent judgment); \emph{Epistle} 63 on the Eucharistic cup.
\item St. Cyril of Jerusalem, \emph{Procatechesis}; \emph{Catechetical Lectures} 3; \emph{Mystagogical Catecheses} 1--5, traditionally attributed to Cyril but sometimes assigned to John of Jerusalem, \sourceurl{https://www.newadvent.org/fathers/3101.htm}{working English collection}; St. Ambrose, \emph{De mysteriis} 1--9, \sourceurl{https://www.newadvent.org/fathers/3405.htm}{working English text}; \emph{De sacramentis} I--VI.
\item St. Gregory Nazianzen, \emph{Oration 40 on Holy Baptism}, especially 3--8, 25, 40--45; St. Basil, \emph{De Spiritu Sancto} 10.24; 12.26--28; 15.35--36; 27.66--68.
\item St. John Chrysostom, \emph{Baptismal Instructions} 2--4; \emph{Homilies on John} 45--47; \emph{Homilies on First Corinthians} 24 and 27; \emph{Homily 82 on Matthew}.
\item St. Augustine, \emph{De doctrina christiana} II.1.1; \emph{Epistle} 98.9; \emph{Tractates on John} 6.7, 26.11--18, 80.3, 120.2; \emph{Sermons} 227, 272; \emph{Contra epistulam Parmeniani} II.13.28; \emph{De baptismo contra Donatistas} I--VII.
\item St. Ignatius of Antioch, \emph{Letter to the Smyrnaeans} 7--8; St. Ephrem the Syrian, \emph{Hymns on Faith} 10 and \emph{Hymns on the Nativity} 3, as Syriac witnesses to Eucharistic fire, Spirit, and embodied mystery.
\end{itemize}
}

\subsection*{Patristic foundations for healing}
{\footnotesize
\begin{itemize}[itemsep=0.18em,topsep=0.25em]
\item \emph{Didache} 4.14; 14.1; \emph{Shepherd of Hermas}, Mandate IV; Tertullian, \emph{De paenitentia} 4--12; St. Cyprian, \emph{De lapsis} 15--36 and selected epistles on reconciliation.
\item Origen, \emph{Homilies on Leviticus} II.4; St. Ambrose, \emph{De paenitentia} I--II; St. John Chrysostom, \emph{On the Priesthood} III.5--6; St. Augustine, \emph{Sermons} 351--352; \emph{Enchiridion} 64--83; \emph{Tractates on John} 121.3--4.
\item St. Leo the Great, \emph{Letter} 168.2 (AD 459), forbidding compulsory public recitation of written sin lists and witnessing secret disclosure to priests; used as disciplinary history, not as a complete modern ritual description.
\item Pope Innocent I, \emph{Epistle 25 to Decentius of Gubbio} 8 (AD 416); ancient Roman and Eastern prayers for blessing the oil of the sick; St. Caesarius of Arles, \emph{Sermon} 13.3; Venerable Bede, \emph{Commentary on James} 5:14--15. Early direct evidence for Anointing is materially smaller and more dispersed than for Baptism and Eucharist.
\item St. Gregory the Great, \emph{Moralia in Iob} and \emph{Pastoral Rule} III on illness, consolation, and differentiated pastoral care; patristic illness texts are used as doctrinal illumination where they do not comment directly on Mark 6 or James 5.
\end{itemize}
}

\subsection*{Patristic foundations for Orders and Matrimony}
{\footnotesize
\begin{itemize}[itemsep=0.18em,topsep=0.25em]
\item St. Clement of Rome, \emph{First Letter to the Corinthians} 40--44; St. Ignatius of Antioch, \emph{Trallians} 2--3; \emph{Magnesians} 2--7; \emph{Smyrnaeans} 8; \emph{Letter to Polycarp} 5.2.
\item \emph{Apostolic Tradition}, ordination prayers in 2--9, cited with the caution that the composite text's attribution, date, and reconstruction remain debated; St. Cyprian, \emph{Epistles} 55, 66--67; St. John Chrysostom, \emph{On the Priesthood} III.4--6; St. Gregory Nazianzen, \emph{Oration} 2.71; St. Gregory the Great, \emph{Pastoral Rule} II--III.
\item Tertullian, \emph{Ad uxorem} II.8; St. John Chrysostom, \emph{Homily 20 on Ephesians}, \sourceurl{https://www.newadvent.org/fathers/230120.htm}{working English text}; \emph{Homily 62 on Matthew}, \sourceurl{https://www.newadvent.org/fathers/200162.htm}{working English text}; \emph{Homily 12 on Colossians}.
\item St. Augustine, \emph{De bono coniugali}, especially 1.1, 6.6, 18.21, 24.32, \sourceurl{https://www.newadvent.org/fathers/1309.htm}{working English text}; \emph{De Genesi ad litteram} IX; \emph{Exposition on Psalm 127(128)}, \sourceurl{https://www.newadvent.org/fathers/1801128.htm}{working English text}; \emph{Tractates on John} 120.2.
\end{itemize}
}

\subsection*{Catholic Churches \latin{sui iuris} and initiation practice}
{\footnotesize
\begin{itemize}[itemsep=0.18em,topsep=0.25em]
\item The controlling universal norms are CCEO 694--697 and 710, \emph{Orientalium Ecclesiarum} 12--14, and the 1996 Eastern liturgical Instruction 50--51. The Latin comparison uses CIC 867, 891, 914 and CCC 1233--1244.
\item Catholic Near East Welfare Association, \sourceurl{https://cnewa.org/eastern-christian-churches/introduction/}{``Eastern Christian Churches: Introduction''} and the linked current profiles, used to check the names, families, and ecclesial identity of the Eastern Catholic Churches; Holy See, \sourceurl{https://www.vatican.va/content/romancuria/en/dicasteri/dicastero-chiese-orientali/profilo.html}{Dicastery for the Eastern Churches profile}.
\item Chaldean Catholic Diocese of St.~Thomas the Apostle U.S.A., \sourceurl{https://chaldeanchurch.org/sacraments/}{``Sacraments,'' Sacraments of Initiation} and \sourceurl{https://chaldeanchurch.org/policies/}{2023 Diocesan Sacramental Policies}: Baptism and Confirmation together, with infant Eucharist locally delayed; unbaptized children of catechetical age receive B--C--E together.
\item Maronite Catholic Eparchy of Saint Maron of Montreal, \sourceurl{https://www.maronites.ca/maronites/maronite-liturgy/}{``Maronite Liturgy,'' Baptism and Confirmation}; Eparchy of Saint Maron of Brooklyn, Beggiani, \sourceurl{https://www.stmaron.org/anaphora}{``The Anaphora,'' Distribution of Communion}; St.~Anthony Maronite Church, \sourceurl{https://www.stanthonylawrence.org/curriculum/}{Eparchy curriculum, Grade 2}. Together these distinguish joined B+C, the historical cessation of infant Communion, and a documented present parish First-Communion timetable; the parish timetable is not universal law.
\item Syro-Malabar Synodal Commission for Liturgy, \sourceurl{https://www.syromalabarliturgy.org/assets/uploads/pdfs/Sacraments-English.pdf}{\emph{The Sacraments of the Syro-Malabar Church}}, General Instructions 14--15 (printed p.~11 / PDF p.~6), directs Communion as early as possible after Baptism and Chrismation; Paul Pallath, \sourceurl{https://console.syromalabarchurch.in/uploads/documents/2023_October-_Correct1.pdf}{``The Liturgical Identity of the Syro-Malabar Church''}, pp.~52--59, supplies historical and implementation qualifications as curial-hosted scholarship, not legislation.
\end{itemize}
}
